\chapter{Inscrição secundária}
\label{chap:pf-secundaria}

\section{O que é?}

Inscrição realizada nos casos em que a/o psicóloga/o é inscrita/o em CRP de outra região, mas precisa desenvolver atividades profissionais em São Paulo por período superior a 90 (noventa) dias por ano. A/o profissional poderá solicitar inscrição secundária no CRP~SP, sem cobrança de anuidade.

\section{Quem pode utilizar esse serviço?}

Psicóloga/o inscrita/o em CRP de outra região que precise trabalhar, por período superior a 90 (noventa) dias por ano, no estado de São Paulo.

\section{Como solicitar o serviço?}

Deve ser solicitado de forma \emph{on-line}, pelo \href{https://www.crpsp.org/pagina/view/338}{\emph{site} do CRP~SP}.

\section{Que informações ou documentos são necessários?}

\begin{itemize}
	\item
	Documento de identificação (RG, CNH etc.);
	\item
	Cadastro de Pessoa Física (CPF);
	\item
	certidão de quitação eleitoral;
	\item
	comprovante ou declaração de endereço;
	\item
	diploma de formação em Psicologia (frente e verso) ou certidão de
	colação de grau;
	\item
	certidão de casamento ou de união estável;
	\item
	certidão de quitação militar ou certificado de reservista (para homens
	até 45 anos);
	\item
	CIP do CRP de origem;
	\item
	declaração com o motivo da solicitação.
\end{itemize}

\section{Quais são as etapas para a realização desse serviço?}

\begin{enumerate}
	\item
	Solicitação via \href{https://www.crpsp.org/agendamento}{\emph{site} do CRP~SP};
	\item
	comparecimento na subsede de referência (confira os endereços \href{https://www.crpsp.org/sede/index}{no \emph{site}}), após
	agendamento, para apresentação dos documentos originais e coleta dos dados biométricos.
\end{enumerate}

\section{Quanto custa?}

\begin{itemize}
	\item
	R\$~70,83 (valor de emissão da CIP em 2024; não há cobrança de anuidade).
\end{itemize}

\section{Quanto tempo leva?}

30 dias após a solicitação no \emph{site}.

\section{Como ter acesso ao atual \emph{status} do serviço?}

No \emph{site} do CRP~SP, acessando a página da/o psicóloga/o com seu \emph{login} e senha, no campo \href{https://cfpsp.brctotal.com/crp06_servicosonline/pgsRequerimento/MeusRequerimentos.aspx}{Meus requerimentos}.

\section{Legislação relacionada}

\begin{itemize}
	\item
	\href{https://atosoficiais.com.br/cfp/resolucao-administrativa-financeira-n-3-2007-institui-a-consolidacao-das-resolucoes-do-conselho-federal-de-psicologia?origin=instituicao&q=resolu\%C3\%A7\%C3\%A3o\%203/2007}{Resolução CFP nº~03/2007} (Consolidação das Resoluções do Conselho Federal de Psicologia);
	\item
	\href{https://transparencia.cfp.org.br/wp-content/uploads/sites/7/2020/05/Manual-de-Procedimentos-Administrativos-e-Financeiro-CFP-RES-20.2018.pdf}{\href{https://atosoficiais.com.br/cfp/resolucao-administrativa-financeira-n-20-2018-altera-o-manual-de-procedimentos-administrativo-e-financeiro-do-sistema-conselhos-de-psicologia-anexo-da-resolucao-cfp-n-202018-a-resolucao-cfp-n-03-2007-a-resolucao-cfp-n-16-2019-e-da-outras-providencias}{Resolução CFP nº~20/2018}}, que determina a revisão e ampliação do
	\href{https://transparencia.cfp.org.br/wp-content/uploads/sites/7/2020/05/Manual-de-Procedimentos-Administrativos-e-Financeiro-CFP-RES-20.2018.pdf}{\emph{Manual de procedimentos administrativos e financeiros}}.
\end{itemize}

\section{Serviços correlatos}

\begin{itemize}
	\item
	\hyperref[chap:pf-inscricao]{Inscrição profissional de pessoa física};
	\item
	\hyperref[chap:pf-cip]{entrega da carteira de identidade profissional (CIP)}.
\end{itemize}

\sumario